\documentclass{article}
\usepackage[margin=0.7in]{geometry}
\usepackage{graphicx}
\usepackage{subcaption}
% ---- preamble: one helper so every cell gets the same size cap ----
\newcommand{\gridimg}[1]{%
  \includegraphics[width=\linewidth,height=0.125\textheight,keepaspectratio]{#1}}

\title{SNN: TODO and test (12th Oct - )}
\author{Omar Abbas Haque}
\date{October 2026}

\begin{document}

\maketitle

\section{What tests are worth doing}

\subsection{Why 9 is biased post-accuracy spike}
Theoretically, if we removed the bias towards 9, we get a near-100 percent accuracy network. Try normalizing the inputs

\subsection{Network designs without an accuracy spike}
Might be worth running the previous test suite over a network design with no accuracy spike

\subsection{Less inhibition}
"200.0 mV inhibition causes one neuron to win the entire window". Does it really? \\
Then try architectures with less inhibition (or tell Sunny to do that)

\section{Observations}
\subsection{50 neurons are doing all the work:}
Looking at F2 for theta\_plus=1 tests, accuracy spike happens at 350 untrained neurons. Actually maybe not, because the highest percentage of spikes come from untrained neurons

\end{document}

